% Part: sets-functions-relations
% Chapter: functions
% Section: functions-relations

\documentclass[../../../include/open-logic-section]{subfiles}


\begin{document}

\olfileid{sfr}{fun}{rel}

\olsection{Functies als relaties}

\begin{explain}
Een functie die !!{element}s van~$A$ afbeeldt op !!{element}s van~$B$
bepaalt vanzelf een relatie tussen $A$ en~$B$, namelijk de relatie die
tussen $x$ en $y$ geldt dan en slechts dan als $f(x) = y$. We kunnen
de functie~$f$ zelfs---bijvoorbeeld als we het aantal bouwstenen van
de wiskunde willen beperken---\emph{identificeren} met deze relatie,
dus met een verzameling paren. Dit roept de vraag op: welke relaties
bepalen op deze manier functies?
\end{explain}

\begin{defn}[Grafiek van een functie] Laat $f\colon A \to B$ een
functie zijn. De \emph{grafiek} van~$f$ is de relatie $R_f \subseteq A
\times B$ die wordt gedefinieerd door
\[
R_f = \Setabs{\tuple{x,y}}{f(x) = y}.
\]
\end{defn}

\begin{explain}
Uit extensionaliteit volgt dat de grafiek van een functie uniek is.
Bovendien rechtvaardigt extensionaliteit voor verzamelingen direct het
stilzwijgende extensionaliteitsbeginsel voor functies: als $f$ en~$g$
hetzelfde domein en codomein hebben, zijn ze gelijk wanneer ze overal
dezelfde waarden aannemen.

Evenzo is een ‘functionele’ relatie de grafiek van een functie.
\end{explain}

\begin{prop}\ollabel{prop:graph-function}
Laat $R \subseteq A \times B$ aan de volgende voorwaarden voldoen:
\begin{enumerate}
\item Als $Rxy$ en $Rxz$, dan $y = z$; en
\item voor iedere $x \in A$ bestaat er een $y \in B$ waarvoor
$\tuple{x, y} \in R$.
\end{enumerate}
Dan is $R$ de grafiek van de functie $f\colon A \to B$ die wordt
gedefinieerd door $f(x) = y$ dan en slechts dan als $Rxy$.
\end{prop}

\begin{proof}
Stel dat er een $y$ bestaat waarvoor $Rxy$ geldt. Als er nog een $z
\neq y$ bestond waarvoor $Rxz$ geldt, zou niet aan de voorwaarde
voor~$R$ zijn voldaan. Als er een $y$ bestaat waarvoor $Rxy$ geldt, is
deze $y$ dus uniek en is $f$ goed gedefinieerd. Uiteraard is $R_f =
R$.
\end{proof}

\begin{explain}
Iedere functie $f\colon A \to B$ heeft een grafiek, dat wil zeggen een
relatie die bestaat uit paren in $A \times B$ en wordt gedefinieerd
door $f(x) = y$.
Omgekeerd is iedere relatie~$R \subseteq A \times B$ met de
eigenschappen uit \olref{prop:graph-function} de grafiek van een
functie~$f \colon A \to B$. Door dit nauwe verband tussen functies en
hun grafieken kunnen we een functie eenvoudig als haar grafiek
opvatten. Functies kunnen met andere woorden worden geïdentificeerd
met bepaalde relaties, dus met bepaalde verzamelingen tupels.
\oliflabeldef{sfr:rel:ref:sec}{Merk echter op dat deze ‘identificatie’
dezelfde strekking heeft als in \olref[sfr][rel][ref]{sec}: het is
geen bewering over de metafysica van functies, maar de vaststelling
dat het handig is functies als bepaalde verzamelingen te
\emph{behandelen}. Een voordeel daarvan is het volgende: w}{W}e
kunnen nu op functies soortgelijke bewerkingen uitvoeren als we op
relaties hebben uitgevoerd (zie \olref[sfr][rel][ops]{sec}). In het
bijzonder:
\end{explain}

\begin{defn}\ollabel{defn:funimage}
Laat $f \colon A \to B$ een functie zijn met $C\subseteq A$.

De \emph{restrictie} van~$f$ tot~$C$ is de
functie~$\funrestrictionto{f}{C}\colon C \to B$ die wordt gedefinieerd
door $(\funrestrictionto{f}{C})(x) = f(x)$ voor alle $x \in C$. Met
andere woorden, $\funrestrictionto{f}{C} = \Setabs{\tuple{x, y} \in
R_f}{x \in C}$.

Als we~$f$ op~$C$ \emph{toepassen}, krijgen we $\funimage{f}{C} =
\Setabs{f(x)}{x \in C}$. Dit noemen we ook het \emph{beeld} van~$C$
onder~$f$.
\end{defn}

\begin{explain}
Uit deze definities volgt dat $\ran{f} =
\funimage{f}{\dom{f}}$ voor iedere functie~$f$.
\oliflabeldef{sfr:rel:ops:sec}{Deze begrippen werken precies zoals men
op grond van de definities in \olref[sfr][rel][ops]{sec} en onze
identificatie van functies met relaties zou verwachten. Twee andere
bewerkingen---inversen en samenstellingen van relaties---vereisen echter iets
meer uitleg. Die geven we in \olref[inv]{sec} en
\olref[cmp]{sec}.}{}
\end{explain}
\end{document}
